
\section{[\ 附录一\ ]常见文言多义词义项举要}

\textbf{1. 去}

①\ 离开：登斯楼也，则有去国怀乡，忧谗畏讥，满目萧然，感极而悲者矣。（《岳阳楼记》）
②\ 距离：当是时，项羽军在鸿门，沛公军在霸上，相去四十里。（《鸿门宴》）

③\ 去掉，除掉：去死肌，杀三虫。（《捕蛇者说》）

④\ 消失：四五日外，色香味尽去矣。（《荔枝图序》）

\textbf{2. 从}

①\ 跟随：赵王遂行，相如从。（《廉颇蔺相如列传》）

②\ 顺从：小惠未遍，民弗从也。（《曹刿论战》）

③\ 依从，听从。公子从其计，请如姬。（《信陵君窃符救赵》）

④\ 随行，侍从，读zòng：忠之属也，可以一战。战则请从。（《曹刿论战》）

⑤\ 堂房亲属，升死，其印为予群从所得，至今保藏。（《活板》）

⑥\ 由，自：从小丘西行百二十步······（《至小丘西小石潭记》）

\textbf{3. 使}

①\ 命令，派遣：秦王使使者告赵王，欲与王为好，会于西河外渑池。（《廉颇蔺相如列传》）

②\ 令，让，叫：使天下之人，不敢言而敢怒。（《阿房宫赋》）

③\ 奉使命出使，读shì：计未定，求人可使报秦者，未得。（《廉颇蔺相如列传》）

④\ 奉使命的人，使臣：会使辙交驰，北邀当国者相见······（《指南录后序》）

⑤\ 假使：使六国各爱其人，则足以拒秦。（《阿房宫赋》）

\textbf{4. 谢}

①\ 道歉：入而徐趋，至而自谢。（《触龙说赵太后》）

②\ 辞别：侯生视公子色终不变，乃谢客就车。（《信陵君窃符救赵》）

③\ 告诉：多谢后世人，戒之慎勿忘！（《孔雀东南飞》）

④\ 道谢：哙拜谢，起，立而饮之。（《鸿门宴》）

\textbf{5. 遗}

①\ 丢失：小学而大遗，吾未见其明也。（《师说》）

②\ 遗留：初患烈遗言：“我死当葬梅花岭上。”（《梅花岭记》）

③\ 遗传下来的：诚宜开张圣听，以光先帝遗德。（《出师表》）

④\ 送，读 wèi：公子姊······数遗魏王及公子书，请教于魏。（《信陵君窃符救赵》）

⑤\ 赠送，读 wèi：人贱物亦鄙，不足迎后人。留待作遗施，于今无会因。（《孔雀东南飞》）

\textbf{6. 策}

①\ 竹制的马鞭：振长策而御宇内。（《过秦论》）

②\ 鞭策、鞭打：策之不以其道，食之不能尽其材。（《马说》）

③\ 骑上鞭打：乃蹶蒙氏衣白衣，复自策其马，麾众拥蒙民马前······（《书博鸡者事》）

④\ 记在简策上：策勋十二转，赏赐百千强。（《木兰诗》）

⑤\ 策划、谋划：请为大王策之。（《左传》）

⑥\ 计策、计谋：均之二策，宁许以负秦曲。（《廉颇蔺相如列传》）

\textbf{7. 间}

①\ 间隙、夹缝：彼节者有间，而刀刃者无厚。（《庖丁解牛》）

②\ 隔离、离间：屈平正道直行······谗人间之，可谓穷矣。（《屈原列传》）

③\ 参与，置身其中；肉食者谋之，又何间焉。（《曹刿论战》）

④\ 间或，断断续续地：数月之后，时时而间进。（《邹忌讽齐王纳谏》）

⑤\ 秘密地，暗地里：又间令吴广之次所旁丛祠中，夜篝火，狐鸣呼曰······（《陈涉世家》）

⑥\ 抄近路，抄小路：从鄱山下，道芷阳间行。（《鸿门宴》）

⑦\ 中间：出师一表真名世，千载谁堪伯仲间！（《书愤》）

⑧\ 一会儿：立有间。（《扁鹊见蔡桓公》）

\textbf{8. 致}

①\ 送给：远方莫不效其珍。（《荀子·解蔽》）

②\ 得到：家贫，无从致书以观。（《送东阳马生序》）

③\ 达到：秦以区区之地，致万乘之势。（《过秦论》）

④\ 表达：一篇之中，三致意焉。（《屈原列传》）

⑤\ 招引：士以此方数千里争往归之，致食客三千人。（《信陵君窃符救赵》）

⑥\ 使至，到：假舆马者，非利足也，而致千里。（《劝学》）

⑦\ 尽，极：衡善机巧，尤致思于天文阴阳历算。（《张衡传》）

\textbf{9. 适}

①\ 到（某地）去：逝将去女，适彼乐土。（《硕鼠》）

②\ 出嫁：贫贱有此女，始适还家门。（《孔雀东南飞》）

③\ 适合、依照：处分适兄意，那得自任专。（《孔雀东南飞》）

④\ 恰好，从上观之，适与地平。（雁荡山）⑤刚才，时已过午，奴辈适至。（《游黄山记》）

\textbf{10.书}

①\ 写：乃丹书帛曰“陈胜王”。（《陈涉世家》）

②\ 文字：读其书未毕，齐军万弩俱发。（史记·吴起孙膑列传）

③\ 书信：忽报曹操遣使送书至（《赤壁之战》）

④\ 文书：昼断狱，夜理书。（《汉书·刑法志》）

⑤\ 书籍：家贫，无从致书以观。（《送东阳马生序》）

⑥\ 特指《尚书》：“《书》曰：“满招损，谦得益。”（《伶官传序》）

\textbf{11. 道}

①\ 道路：会天大雨，道不通。（《陈涉世家》）

②\ 途径、方法：此五者，知胜之道也。（《谋攻》）

③\ 道理：虽有至道，弗学不知其善也。（《学记》）

④\ 规律：臣之所好者，道也；进乎技矣。（《庖丁解牛》）

⑤\ 学说、思想：既加冠，益慕圣贤之道。（《送东阳马生序》）

⑥\ 述说：此中人语云：“不足为外人道也。”（《桃花源记》）

\textbf{12. 迁}

①\ 迁移、迁都：时北兵已迫修门外，战、守、迁皆不及施。（《指南录后序》）

②\ 变易：齐人未尝赂秦，终继五国迁灭，何哉？（《六国论》）

③\ 调动、升官：公车特征，拜郎中，再迁为太史令。（《张衡传》）

④\ 降职，贬谪；迁客骚人，多会于此。（《岳阳楼记》）

⑤\ 放逐；顷襄王怒而迁之。（《屈原列传》）

\textbf{13. 卒}

①\ 步兵；夫以疲病之卒，御狐疑之众······（《赤壁之战》）

②\ 一百人的队伍；凡用兵之法······全卒为上，破卒次之。（《谋攻》）

③\ 完毕；语未及卒，公子立变色。（《信陵君窃符救赵》）

④\ 终于；卒相与欢，为刎颈之交。（《廉颇蔺相如列传》）

⑤\ 死；初，鲁肃闻刘表卒。（《赤壁之战》）

⑥\ 突然（在这个意义上，与“猝”通，读cù）；五万兵难卒合，已选三万人······（《赤壁之战》）

\textbf{14. 举}
①\ 举起、抬起；千钧之重，人不轻举。（《盐铁论·刑德》）

②\ 升起；中江举帆，余船以次俱进。（《赤壁之战》）

③\ 举荐；举孝廉不行，连辟公府不就。（《张衡传》）

④\ 发动；今亡亦死，举大计亦死，等死，死国可乎？（《陈涉世家》）

⑤\ 攻下；戍卒叫，函谷举。（《阿房宫赋》）

⑥\ 全，尽；杀人如不能举。（《鸿门宴》）

⑦\ 总合；吾不能举全吴之地，十万之众，受制于人。（《赤壁之战》）

\textbf{15. 负}

①\ 依仗、负恃；秦负，负其强，以空言求璧。（《廉颇蔺相如列传》

②\ 背靠着：负海之国贵攻战。（《商君书·兵守》）

③\ 用背驮东西，背着：当余之从师也，负箧曳履，行深山巨谷中······（《送东阳马生序》）

④\ 担负：均之二策，宁许以负秦曲。（《廉颇蔺相如列传》）

⑤\ 背弃：相如度秦王虽斋，决负约不偿城。（《廉颇蔺相如列传》）

⑥\ 辜负：臣诚恐见欺于王而负赵。（《廉颇蔺相如列传》）

⑦\ 失败：不知彼而知己，一胜一负。（《谋攻》）

为了节省篇幅，亦为了让读者自己有意识地到语文课本中去寻找包含如下一些常见文言多义词各个义项的句子，下面例句一律从略。

\begin{enumerate}[label=\textbf{\arabic*. }, start=16]
	\item 及 --- 追赶上；达到；涉及；趁着；比得上；与。
	\item 相 --- 仔细看（或“审察”）；容貌；辅助；扶助君王的人（宰相）；互相；表示动作偏指一方。
	\item 奉 --- 两手恭敬地捧着；进献；恭敬地接受；遵从；拥戴；供奉；俸禄（这个意义后来写作“俸”）。
	\item 素 --- 没有染色的丝绸；白的；空；朴素；向来。
	\item 市 --- 市场；买；交易（常用于比喻义）。
	\item 顾 --- 回头看；看望；关照；思念；不过。
	\item 游 --- 在水上漂浮；浮动；虚浮；游玩；出外求学或求官；交往。
	\item 坐 --- 两膝靠着席子而且臀部压在脚根上的姿势；座位（这个意义后来写为“座”）；犯罪；因为。
	\item 士 --- 男子（特指未婚男子），战车上的兵士，统治阶级的下层；读书人，具有某种品质或某种技能的人。
	\item 令 --- 命令，使，假使，美好，县令，时令。
	\item 疾 --- 病，毛病，痛苦，痛恨，急速，强烈妒忌（这个意义又写为“嫉”）。
	\item 集 --- 群鸟停在树上，停留，聚合，诗文的汇集。
	\item 临 --- 从高处往低处看，从上监视着（或“统治”）；到，面对，对着书画范本摹仿学习。
	\item 食 --- 吃，吃的东西，供养（或者“给······吃”，这个意义读 sī），饲养（也读作 sì）。
	\item 许 --- 应允，赞同，表示大约的数量，处所（常与疑问代词“何”、“恶”等连用）。
	\item 居 --- 坐，处在，居住，住处，占据，表示相隔了一段时间（用在“有顷”，“久之”，“顷之”等前面）。
	\item 绝 --- 绳索断，断，隔绝，到了极点的，高超，横渡。
	\item 然 --- 燃烧（这个意义后来写作“燃”）；这样（或“那样”）；是的（或“对”）；认为······是对的，······的样子（······地）。
	\item 故 --- 原因，事故，旧，原来的，旧交，故意，特意，因此。
	\item 上 --- 位置在高处的，地位高的，帝王，质量高的，次序在前的，登上，进献。
	\item 师 --- 军队二千五百人，军队；老师；学习；乐官。
	\item 说 --- 解说；言论（或“主张”），说服（或“劝说”，读 shuì）；高兴（这个意义后来写为“悦”）。
	\item 徒 --- 步行；步兵；徒党（同一类的人）；门徒（徒弟）；空；白白地；只（仅仅）。
	\item 事 --- 事情；从事（做）；侍奉（为······服务）；对待。
	\item 数 --- 数目（数量）；几（几个）；计算；一列举（上面这两项都读 shù）；屡次（读 shuò）；密（读 cù）；命运；必然性；技艺。
	\item 固 --- 坚固；坚持；顽固；本来；当然；一定。
	\item 器 --- 器具；才能；器重；气量。
	\item 度 --- 量长短；估量；衡量；（读 dù）量长短的标准（读 dù）；限度；尺度；制度；法度；气度；渡过（这个意义后来又写为“渡”）。
	\item 患 --- 担忧；忧患；灾祸；害疾病；得······病。
	\item 穷 --- 阻塞不通；境遇不好；不能显贵；生活困难；尽头；寻究。
	\item 殆 --- 危险；近于；几乎，恐怕；也许。
	\item 方 --- 两船平行；比拟；方（跟“圆”相对）；正直；范围；方向；方法；正在。
	\item 会 --- 会合；盟会；会见；机会；正巧；一定；算帐（读 kuài）。
	\item 让 --- 责备；辞让；退让；谦让；让权位。
	\item 攘 --- 排斥；偷；侵夺；撩起；挽起。
	\item 征 --- 远行；征伐；征税；征召；证明；征兆。
	\item 责 --- 索取；要求；责令；责备；责任；欠别人的钱财（这个意义后来写作“债”，读zhài）。
	\item 要 --- 腰（这个意义后来写作“腰”），读yāo）；半路拦截，邀请；求得；要挟（以上各词义均读yāo）；重要的；总括；需要。
	\item 伐 --- 砍伐；讨伐；声讨；攻打；攻破；夸耀。
	\item 属 --- 连接；跟随；撰写（“属文”的“属”）；委托；嘱托（这个意义后来写作“嘱”）。（按：以上各词义均读zhǔ）。隶属；掌管；部属；亲属；类（一类的人）。（按：以上各词义均读shù）。
	\item 安 --- 安定；安乐；安心；安放；什么（或“什么地方”）；怎么（或“哪里”）。
	\item 暴 --- 晒（这个意义后来又写作“曝”，读pù）；暴露；又猛又急；突然；凶恶残酷的；欺凌。
	\item 被 --- 被子；覆盖；加于······之上；遭受；蒙受；表示被动；穿在身上（或“披在身上”）。这个意义后来写作“披”，读pī）。
	\item 或 --- 有的；有的人；有时；也许；用在否定句中加强否定语气。
	\item 是 --- 正；正确；这；这儿；这样；这个。
\end{enumerate}
\newpage